\section{Coverage ledger and source-by-source audit}\label{sec:coverage-audit}

This section records exactly what has been proved, in which sense, and where it
occurs.  The letters in the last column use the four status levels of
\cref{sec:scope-status}: \statusA{} is an exact arithmetic theorem already
represented by executable finite code; \statusB{} is a fully quantified
finite contract with explicit moduli; \statusC{} is a mathematically complete
constructive reduction whose routine certificate compiler remains to be
written; and \statusD{} is deliberately outside the mathematical target or is
retained only as a source annotation.  No row labelled \statusA, \statusB, or
\statusC{} is meant to assert that the final object-language derivation
\( []\vdash_{\mathrm{PRAAx}}\varphi\) has already been checked by Lean.

\renewcommand{\arraystretch}{1.17}
\small
\begin{longtable}{>{\raggedright\arraybackslash}p{0.07\textwidth}
                  >{\raggedright\arraybackslash}p{0.28\textwidth}
                  >{\raggedright\arraybackslash}p{0.40\textwidth}
                  >{\centering\arraybackslash}p{0.08\textwidth}}
\toprule
ID & Source statement & Constructive proof and qualification & Status\\
\midrule
\endfirsthead
\toprule
ID & Source statement & Constructive proof and qualification & Status\\
\midrule
\endhead
\midrule
\multicolumn{4}{r}{Continued on next page}\\
\endfoot
\bottomrule
\endlastfoot

W1 & Weierstrass canonical product for \(G(1+z)\) &
The product is defined by finite logarithmic factors.  The estimate
\(\lvert q_k(z)\rvert\le 2R^3/(3k^2)\), valid for
\(\lvert z\rvert\le R\) and \(k\ge2R\), gives a primitive-recursive cutoff,
uniform Cauchy modulus, and product enclosure; see \cref{lem:canonical-tail}.
& \statusB\\

W2 & Improper integral representation &
Direct convergence is proved for \(\Re z>-1\).  The interval is split at a
primitive-recursive truncation point; Taylor bounds near zero, exponential
tails at infinity, and rational quadrature on the middle interval give a finite
certificate.  Outside this half-plane one must use a separately certified
analytic continuation path; see \cref{thm:improper}. & \statusB\\

W3 & Entire function &
Uniform convergence of the canonical product and all fixed derivatives on
rational compact disks is reduced to explicit \(p\)-series tails.  The
constructive statement is a compact-disk family, not an appeal to a completed
space of holomorphic functions; see \cref{thm:entire}. & \statusB\\

W4 & Order two and infinite type &
The upper bound follows from the canonical product by splitting at
\(k\le2r\); the lower bound uses exact integer values and elementary factorial
bounds.  The resulting finite inequalities imply order two and infinite type;
see \cref{thm:order-type}. & \statusB\\

W5 & Functional equation \(G(z+1)=\Gamma(z)G(z)\) &
Finite products cancel exactly.  Harmonic and Stirling remainders identify the
remaining finite quotient with the Gamma Euler product.  Equality is exported
as cross-multiplied complex-ball inclusions on apartness-certified domains;
see \cref{thm:functional}. & \statusB\\

W6 & Normalization \(G(1)=1\) &
The defining product has its zero argument and every factor equal to one.  The
claim is exact and needs no limit. & \statusA\\

W7 & Integer values and superfactorials &
Quantifier-free induction on the functional equation gives
\(G(n+1)=\prod_{j<n}j!\).  The same recurrence proves the superfactorial
identification and all displayed initial values; see
\cref{thm:integer-ipra}. & \statusA\\

W8 & Hyperfactorial/K-function relation &
Both sides are reduced to finite products at integers and to the common
Gamma--Barnes recurrence for named positive-real arguments.  Division is
cleared before export to iPRA; see \cref{sec:hurwitz}. & \statusB\\

W9 & Vign\'eras third-derivative characterization &
The canonical product gives
\((\log G)'''(x)=2\sum_{k\ge0}(k+1)/(x+k)^3>0\), with an explicit tail.
For a competitor, the logarithmic quotient is periodic; positivity at large
integer translates forces its third derivative to be nonnegative, while its
period integral is zero.  Hence it is constant; see \cref{thm:vigneras}.
The coded-function compiler is still a formalization task. & \statusC\\

W10 & Duplication formula &
This is the case \(n=2\) of the multiplication theorem.  All constants are
obtained by exact finite rearrangement and the value at one half; see
\cref{sec:multiplication}. & \statusB\\

W11 & Asymptotic characterization &
The source display is corrected to
\(f(n+x)\sim\Gamma(n)^x n^{\binom{x}{2}}f(n)\).
A uniform compact-\(x\) remainder follows from the digamma strip estimate and
Euler--Maclaurin.  Periodicity of the quotient then gives uniqueness; see
\cref{thm:corrected-characterization}. & \statusB\\

W12 & Kinkelin reflection formula &
Differentiate both sides on a rational path avoiding the integers, use the
Gamma reflection identity, and fix the constant at the base point.  The
proof is formulated as path integration with a zero-free certificate; see
\cref{thm:kinkelin}. & \statusB\\

W13 & Clausen reflection forms &
Insert the finite sine-series approximation to the Clausen function in the
reflection formula.  A geometric tail bound and rational trigonometric balls
supply the modulus; see \cref{sec:reflection}. & \statusB\\

W14 & Half-shift tangent identity &
Subtract the reflection formulas at two translated arguments and use the
functional equation and elementary trigonometric identities.  The domain is a
compact interval quantitatively separated from half-integers. & \statusB\\

W15 & Taylor series for \(\log G(1+z)\) &
Termwise expansion of each finite canonical factor is justified first, then a
uniform tail permits passage to the limit.  For \(\lvert z\rvert\le\rho<1\),
the tail after degree \(M\) is at most
\(2\rho^{M+1}/((M+1)(1-\rho))\); see \cref{thm:taylor}. & \statusB\\

W16 & Exponentiated Taylor/product relation &
A rational complex exponential polynomial with its explicit tail converts the
logarithmic enclosure into an enclosure for \(G\).  No global logarithm is
asserted. & \statusB\\

W17 & General multiplication formula &
A finite double product is regrouped by residue classes modulo \(n\).  The
remaining quotient is periodic; the asymptotic formula fixes its constant.
The result is exported in exponentiated form with branch data; see
\cref{thm:multiplication}. & \statusB\\

W18 & Multiplication constant &
The constant is reduced to powers of \(2\pi\), \(n\), and the named Glaisher
constant.  For each fixed \(n\), all integer exponents and products are
primitive recursive; see \cref{sec:multiplication}. & \statusB\\

W19 & Conjugation &
Every rational finite product commutes exactly with conjugation, and the same
cutoff works for \(z\) and \(\bar z\). & \statusA\\

W20 & Absolute-value product &
Multiply the product at \(z\) by its conjugate.  The preferred statement is
for \(\lvert G(z)\rvert^2\), avoiding square roots.  Positivity and a lower
bound permit extraction of the unsquared modulus when needed. & \statusB\\

W21 & Pure-imaginary modulus &
The source prefactor \(y\) is corrected to \(\lvert y\rvert\).  The squared
identity has the exact factor \(y^2\) and is valid without a sign decision; see
\cref{audit:modulus}. & \statusB\\

W22 & Barnes asymptotic expansion &
Euler--Maclaurin with the periodic Bernoulli remainder yields the target
series.  The Nemes bound supplies a closed error term for every fixed order and
sector; see \cref{thm:Nemes-bound}. & \statusB\\

W23 & Log-Gamma integral &
Alexeiewsky's identity is proved first for finite quadrature data by
integration by parts and the functional equations, then by explicit mesh and
tail estimates; see \cref{thm:alex}. & \statusB\\

W24 & Logarithmic derivative and derivative formulas &
Differentiate the canonical logarithm on an apartness-certified disk.  The
quotient form is replaced by
\(G'=G\cdot (G'/G)\) when no lower bound for \(G\) is available. & \statusB\\

M1 & Integer hyperfactorial formula &
Finite products and the recurrence give the identity by quantifier-free
induction. & \statusA\\

M2 & Initial integer sequence &
The first values are numeral normalizations of the primitive-recursive
superfactorial.  The companion checker verifies a wider range. & \statusA\\

M3 & Hurwitz-zeta derivative formula &
Euler--Maclaurin defines a fast name for \(\zeta'(s,a)\) at \(s=-1\), with an
explicit remainder.  Subtracting the same recurrence on both sides and fixing
the value at \(a=1\) proves Vardi's formula; see \cref{thm:vardi}. & \statusB\\

M4 & Negative-order polygamma notation &
The notation is made definitional rather than imported ambiguously: the
chosen normalized antiderivative is tied to the Hurwitz-zeta derivative and a
base value.  Under this normalization the displayed Barnes formula follows
from M3. & \statusB\\

M5 & Finite Gamma product &
Repeated use of the functional equation gives a division-free product identity
with exact finite induction. & \statusA\\

M6 & Derivative formula &
This is W24 rewritten using Alexeiewsky's identity and the digamma fast name.
All divisions require an apartness witness. & \statusB\\

M7 & Leading Stirling-like asymptotic &
This is the first truncation of W22, with the remainder specialized to a closed
bound. & \statusB\\

M8 & Full Bernoulli asymptotic &
For every externally fixed truncation order, the coefficients and the Nemes
remainder bound are primitive-recursive names.  Uniformity is stated only on a
fixed rational sector. & \statusB\\

M9 & \(G(1/4)\) &
Duplication, reflection, the Gamma quarter reflection formula, and the named
Catalan and Glaisher constants determine the positive value.  Branches are
fixed by positivity. & \statusB\\

M10 & \(G(3/4)\) &
The same system determines the value; the two MathWorld forms are equivalent by
\(\Gamma(1/4)\Gamma(3/4)=\pi\sqrt2\). & \statusB\\

M11 & \(G(1/2)\) &
The duplication constant and normalization give the standard formula, with
positive-real branches. & \statusB\\

M12--M13 & \(G(3/2)\), \(G(5/2)\) &
Apply the functional equation to M11 and use the exact Gamma half-integer
values. & \statusB\\

M14 & All odd half-integer values and \(c_k\) &
Induction isolates an exact integer/rational coefficient \(c_k\) multiplying
the fixed named constants.  The recurrence for \(c_k\) is primitive recursive;
see \cref{thm:odd-half}. & \statusA/\statusB\\
\end{longtable}
\normalsize

\subsection{Items on the pages that are not additional mathematical theorems}

The historical notes, software syntax, graphs, decimal evaluations, and lists
of applications are not independent identities to prove in iPRA.  A decimal
printout is covered by a stronger statement here: a rational interval with a
checked width and containment certificate.  Likewise, an implementation note
about a computer algebra system is replaced by the explicit finite algorithms
in the archive.  The report does not attempt to formalize bibliographic claims
such as priority or naming conventions.

The asymptotic Stokes analysis beyond a fixed closed sector is also not needed
to establish the formulas on the two target pages.  It would require an
additional exponentially improved expansion and branch-dependent terminant
bounds.  Nothing in this report uses such a claim silently.

\subsection{What the ledger establishes}

The ledger covers every substantive displayed identity or characterization on
the two target pages, after the two necessary corrections.  Exact finite
identities are reduced to primitive recursion and quantifier-free induction;
analytic identities are reduced to finite rational computations with explicit
cutoffs, separation hypotheses, and error inequalities.  The remaining work is
therefore not discovery of an unspecified constructive foundation: it is the
mechanical Lean implementation of the indicated contracts and their export to
ProveIt's \code{Prf PRAAx} calculus.
